% HQIV Lean: combinatorial backbone -> electroweak closure -> fermion/hadron layers.
%
% SUPERSEDED — archived 2026-05-31. Content merged into
% papers/tuft_hqiv_dynamic_mass_spectrum/hqiv_tuft_dynamic_mass_spectrum.tex
% (bib key hqiv-tuft-dynamic-mass-spectrum-paper). Do not publish this draft.
%
% pdflatex
\documentclass[11pt]{article}

\usepackage{amsmath,amssymb,amsthm,mathtools}
\usepackage{enumitem}
\usepackage[margin=1in]{geometry}
\usepackage{hyperref}
\hypersetup{colorlinks=true,linkcolor=blue,citecolor=blue,urlcolor=blue}
\usepackage[final]{microtype}
\setlength{\emergencystretch}{3em}

\newcommand{\leanrepo}{https://github.com/HQIV/hqiv-lean/blob/main}
% Lean module / identifier helpers: visible text is rendered with \nolinkurl
% so that slashes, underscores, and dots become legal break points.
\newcommand{\leanfile}[1]{\href{\leanrepo/#1}{\nolinkurl{#1}}}
\newcommand{\leanid}[1]{\nolinkurl{#1}}

\newtheorem{theorem}{Theorem}[section]
\newtheorem{lemma}[theorem]{Lemma}
\newtheorem{proposition}[theorem]{Proposition}
\newtheorem{corollary}[theorem]{Corollary}
\newtheorem{definition}[theorem]{Definition}
\newtheorem{remark}[theorem]{Remark}

\title{Electroweak Scale and Fermion Masses from Horizon Readouts:\\[0.35em]
First-Principles Electroweak Closure inside the HQIV Lean Library}
\author{Steven Ettinger Jr\\
{\small Independent Researcher --- \texttt{steven@disregardfiat.tech}}}
\date{Preprint v1 --- \today}

\begin{document}
\maketitle

\begin{abstract}
\noindent\textbf{Position in the publication chain.}
This is the first paper of the Tier-2 Standard-Model layer (\#5 in the
publication sequence). Tier~0 fixes the framework (HQIV main
paper~\cite{ettinger2026hqiv}, Tier-0a) and the discrete null-lattice readout
(lightcone paper~\cite{hqiv-lightcone-paper}, Tier-0b; closure
paper~\cite{hqiv-so8-paper}, Tier-0c). Tier~1 wraps that spine with
dimensional rigidity \cite{hqiv-3d-growth-paper}, the variational layer
\cite{hqiv-oct-action-paper}, finite-mode Kirchhoff thermodynamics
\cite{hqiv-kirchhoff-paper}, and ladder thermodynamics with the discrete
arrow \cite{hqiv-thermo-arrow-paper}. The present paper rides every one of
those upstream results to derive an electroweak scale and a fermion mass
ladder from horizon readouts alone---no continuum quantum field theory
imported, no Standard-Model Lagrangian assumed in advance. Tier~0 records and
all four Tier~1 companions cited above are published on Zenodo.

\noindent\textbf{Result.}
From the discrete null-lattice combinatorics and horizon-area
readouts already proved upstream, the HQIV Lean formalisation
produces a closed electroweak sector \emph{without importing PDG
inputs for it}. A geometric vacuum
\[
  v=T_{\mathrm{lockin}}\,S(M{+}1)\,(1+\gamma),
  \qquad
  S(m)=(m+1)(m+2),
  \qquad
  M:=\mathrm{referenceM}=4,
\]
is built from the lock-in temperature
$T_{\mathrm{lockin}}=T(M)$ and the outer-horizon surfaces
$S(m)$; the closure paper's lattice $\alpha$-ratio identity
forces $\alpha=3/5$ at every truncation, hence
$\gamma=1-\alpha=2/5$. Gauge-boson and Higgs masses follow from
$v$ with triality-fixed couplings; rational witnesses (for
example $M_W^{\mathrm{derived}}=392/5$ in the module's natural
units, plus matching $M_Z$ and $m_H$) are sorry-free Lean
theorems, and comparison theorems against electroweak reference
values \cite{pdg2024} are also formalised. The same backbone
delivers a derived neutrino ladder, the strong-colour chart on
the $\mathrm{EuclideanSpace}\,\mathbb{C}\,(\mathrm{Fin}\,8)$
carrier shared with the electroweak layer, and a charged-lepton
chain with the lock-in mode-frequency horizon as the primary
input. Reader-friendly anchors with clickable Lean links live in
Appendix~\ref{app:lean-index}.

\noindent\textbf{Single active scale witness.}
The HQIV pipeline carries exactly one dimensionful scale
witness per solve. The default is the
\hyperref[lean:proton-lockin]{proton lock-in scale-witness}, the
same convention that pins the mass/unit chart in the
finite-mode Kirchhoff paper \cite{hqiv-kirchhoff-paper}: the
proton mass at the lock-in shell uses the theory's internal
constituent-plus-binding readout, with the observed value
$M_p=938.272\,\mathrm{MeV}$ \cite{pdg2024} serving as the active
witness. Under that witness, the CODATA $1/\alpha$
\cite{codata2018}, CMB horizons, and PDG quark/lepton centrals
are \emph{predictions or comparison layers}, never simultaneous
solve anchors.

\noindent\textbf{Patch-ontology reader contract.}
Per the patch-theory reader contract imported below, the
discrete patch layer is the ontology. Continuum forms appearing
in this paper---an
$\mathrm{EuclideanSpace}\,\mathbb{C}\,(\mathrm{Fin}\,8)$ carrier
treated as a finite-dimensional Hilbert space, an
$\mathfrak{so}(8)$ Lie-algebra layer used as bookkeeping for
representation labels, the typed weak/Higgs scaffold built on
the proved O-Maxwell action---are de-facto continuum objects
built from finite, discrete data. No continuum-spacetime limit,
no continuum-QFT recovery, and no mesh-refinement programme is
invoked, required, or desired anywhere in this chain.

\noindent\textbf{Status.}
The Lean development is sorry-free against \texttt{mathlib} for
the electroweak closure witnesses, the neutrino ladder, the
informational-energy mass row, the strong-colour chart on the
shared carrier, and the proved reduction identities of the
weak/Higgs variational extension back to the abelian O-Maxwell
core. The quark ladder still carries explicit GeV anchors and
$\mathbb{N}$ shell coordinates; the honest-accounting block in
Section~\ref{sec:quarks} marks each remaining numeral and
explains what is and is not derived from the
$(T_{\mathrm{lockin}},S,\gamma,\alpha,M)$ spine.
\end{abstract}

% Shared reader contract: patch QFT + discrete numerics.
% From papers/*.tex:     % Shared reader contract: patch QFT + discrete numerics.
% From papers/*.tex:     % Shared reader contract: patch QFT + discrete numerics.
% From papers/*.tex:     \input{include/patch_theory_messaging.tex}
% From papers/paper/*.tex: \input{../include/patch_theory_messaging.tex}

\paragraph{Patch QFT and discrete numerics (reader contract).}
HQIV (Horizon-Quantized Informational Vacuum) is formulated as a \emph{patch theory}: quantum fields, actions, and physical readouts are attached to discrete patches---shell indices $m\in\mathbb{N}$, finite spacetime index types (e.g.\ $\mathrm{Fin}\,4$), and horizon-limited charts---rather than to a hypothetical fundamental smooth spacetime obtained as a mesh-refinement or ultraviolet limit.
Accordingly, when this text refers to ``QFT,'' it means \emph{patch QFT}: patching rules, finite measurement ledgers, and variational identities on the discrete spine; it is not the claim that the theory is \emph{defined} by recovering ordinary continuum quantum field theory through such a limit.
The numbers that admit the sharpest statements---mass and coupling ladders, curvature ratios, shell thermodynamics, and rational witnesses in \texttt{HQIV\_LEAN}---are objects of \emph{discrete mathematics} on $\mathbb{N}$ (combinatorial growth, cumulative simplex ledgers) together with the ordered field $\mathbb{R}$ as used in the proof assistant for analysis, bounds, and phenomenological display.
Smooth-manifold or continuum field notation, when it appears, is a \emph{translation layer} for comparison with standard physics literature; it does not supersede the discrete definitions.

\paragraph{A continuum is not required, and in places would destroy these results.}
The discrete patch layer is \emph{not} a numerical approximation to a more fundamental continuum theory awaiting future construction; it is the ontology.  Where the literature treats ``the continuum limit'' as the goal, HQIV treats it as a \emph{calculation approximation} for comparison with standard formulas.  The continuum form is an optional convenience for comparison; the proved discrete statement is what carries the physics.  In several load-bearing places---the curvature imprint $\delta_E(m)$ at shell $m$, the harmonic ladder masses $m_\tau\to m_\mu\to m_e$, the lock-in vev $v=\sqrt{\eta_{\rm paper}\,\Omega_k(m_{\rm lockin};m_{\rm lockin})}$, the proton anchor at $m=\mathrm{referenceM}=4$, the baryon-asymmetry $\eta\propto 6^7\sqrt{3}/(m+1)$, and the rational coefficients $\alpha=3/5$, $\gamma=2/5$, $\alpha_{\rm GUT}=1/42$---taking a literal $m\to\infty$ or sub-Planck continuum limit would \emph{erase} the discrete index that is doing the predictive work.  Continuum forms of these objects are therefore not preferred refinements; they are degenerate, information-losing readouts of the same patch data.

\paragraph{An exception: $\mathfrak{so}(8)$ as a de-facto continuum algebra from a discrete construction.}
Principled exception to the discrete-only stance: the closed Lie algebra $\mathfrak{so}(8) \supseteq G_2\cup\{\Delta\}$ generated from the Fano octonion left-multiplication table and the phase-lift generator $\Delta\in(\mathrm{e}_1,\mathrm{e}_7)$.  Although $\mathfrak{so}(8)$ is a continuous (28-dimensional real) Lie algebra, it is built here entirely from \emph{discrete} ingredients (the seven imaginary octonion units, a finite Fano table, and one $\pi/2$ rotation), and its closure to 28 dimensions is a finite linear-algebra certificate (see the \texttt{Hqiv.SO8Closure} / \texttt{Hqiv.SO8ClosureInterface} stack).  The Lie-group structure is therefore a \emph{de-facto continuum algebra obtained from a discrete construction}: continuous in parameter space, but with no continuum spacetime, no continuum measure, and no UV limit attached.  When this manuscript or its companions write $\exp(t\,\Delta)\in\mathrm{SO}(8)$ or treat $\mathfrak{so}(8)$ rotations as smooth, those are statements internal to a finite-dimensional Lie group whose generators are certified by discrete data; they are \emph{not} a continuum-spacetime commitment and do not require a continuum-QFT scaffolding.

% From papers/paper/*.tex: % Shared reader contract: patch QFT + discrete numerics.
% From papers/*.tex:     \input{include/patch_theory_messaging.tex}
% From papers/paper/*.tex: \input{../include/patch_theory_messaging.tex}

\paragraph{Patch QFT and discrete numerics (reader contract).}
HQIV (Horizon-Quantized Informational Vacuum) is formulated as a \emph{patch theory}: quantum fields, actions, and physical readouts are attached to discrete patches---shell indices $m\in\mathbb{N}$, finite spacetime index types (e.g.\ $\mathrm{Fin}\,4$), and horizon-limited charts---rather than to a hypothetical fundamental smooth spacetime obtained as a mesh-refinement or ultraviolet limit.
Accordingly, when this text refers to ``QFT,'' it means \emph{patch QFT}: patching rules, finite measurement ledgers, and variational identities on the discrete spine; it is not the claim that the theory is \emph{defined} by recovering ordinary continuum quantum field theory through such a limit.
The numbers that admit the sharpest statements---mass and coupling ladders, curvature ratios, shell thermodynamics, and rational witnesses in \texttt{HQIV\_LEAN}---are objects of \emph{discrete mathematics} on $\mathbb{N}$ (combinatorial growth, cumulative simplex ledgers) together with the ordered field $\mathbb{R}$ as used in the proof assistant for analysis, bounds, and phenomenological display.
Smooth-manifold or continuum field notation, when it appears, is a \emph{translation layer} for comparison with standard physics literature; it does not supersede the discrete definitions.

\paragraph{A continuum is not required, and in places would destroy these results.}
The discrete patch layer is \emph{not} a numerical approximation to a more fundamental continuum theory awaiting future construction; it is the ontology.  Where the literature treats ``the continuum limit'' as the goal, HQIV treats it as a \emph{calculation approximation} for comparison with standard formulas.  The continuum form is an optional convenience for comparison; the proved discrete statement is what carries the physics.  In several load-bearing places---the curvature imprint $\delta_E(m)$ at shell $m$, the harmonic ladder masses $m_\tau\to m_\mu\to m_e$, the lock-in vev $v=\sqrt{\eta_{\rm paper}\,\Omega_k(m_{\rm lockin};m_{\rm lockin})}$, the proton anchor at $m=\mathrm{referenceM}=4$, the baryon-asymmetry $\eta\propto 6^7\sqrt{3}/(m+1)$, and the rational coefficients $\alpha=3/5$, $\gamma=2/5$, $\alpha_{\rm GUT}=1/42$---taking a literal $m\to\infty$ or sub-Planck continuum limit would \emph{erase} the discrete index that is doing the predictive work.  Continuum forms of these objects are therefore not preferred refinements; they are degenerate, information-losing readouts of the same patch data.

\paragraph{An exception: $\mathfrak{so}(8)$ as a de-facto continuum algebra from a discrete construction.}
Principled exception to the discrete-only stance: the closed Lie algebra $\mathfrak{so}(8) \supseteq G_2\cup\{\Delta\}$ generated from the Fano octonion left-multiplication table and the phase-lift generator $\Delta\in(\mathrm{e}_1,\mathrm{e}_7)$.  Although $\mathfrak{so}(8)$ is a continuous (28-dimensional real) Lie algebra, it is built here entirely from \emph{discrete} ingredients (the seven imaginary octonion units, a finite Fano table, and one $\pi/2$ rotation), and its closure to 28 dimensions is a finite linear-algebra certificate (see the \texttt{Hqiv.SO8Closure} / \texttt{Hqiv.SO8ClosureInterface} stack).  The Lie-group structure is therefore a \emph{de-facto continuum algebra obtained from a discrete construction}: continuous in parameter space, but with no continuum spacetime, no continuum measure, and no UV limit attached.  When this manuscript or its companions write $\exp(t\,\Delta)\in\mathrm{SO}(8)$ or treat $\mathfrak{so}(8)$ rotations as smooth, those are statements internal to a finite-dimensional Lie group whose generators are certified by discrete data; they are \emph{not} a continuum-spacetime commitment and do not require a continuum-QFT scaffolding.


\paragraph{Patch QFT and discrete numerics (reader contract).}
HQIV (Horizon-Quantized Informational Vacuum) is formulated as a \emph{patch theory}: quantum fields, actions, and physical readouts are attached to discrete patches---shell indices $m\in\mathbb{N}$, finite spacetime index types (e.g.\ $\mathrm{Fin}\,4$), and horizon-limited charts---rather than to a hypothetical fundamental smooth spacetime obtained as a mesh-refinement or ultraviolet limit.
Accordingly, when this text refers to ``QFT,'' it means \emph{patch QFT}: patching rules, finite measurement ledgers, and variational identities on the discrete spine; it is not the claim that the theory is \emph{defined} by recovering ordinary continuum quantum field theory through such a limit.
The numbers that admit the sharpest statements---mass and coupling ladders, curvature ratios, shell thermodynamics, and rational witnesses in \texttt{HQIV\_LEAN}---are objects of \emph{discrete mathematics} on $\mathbb{N}$ (combinatorial growth, cumulative simplex ledgers) together with the ordered field $\mathbb{R}$ as used in the proof assistant for analysis, bounds, and phenomenological display.
Smooth-manifold or continuum field notation, when it appears, is a \emph{translation layer} for comparison with standard physics literature; it does not supersede the discrete definitions.

\paragraph{A continuum is not required, and in places would destroy these results.}
The discrete patch layer is \emph{not} a numerical approximation to a more fundamental continuum theory awaiting future construction; it is the ontology.  Where the literature treats ``the continuum limit'' as the goal, HQIV treats it as a \emph{calculation approximation} for comparison with standard formulas.  The continuum form is an optional convenience for comparison; the proved discrete statement is what carries the physics.  In several load-bearing places---the curvature imprint $\delta_E(m)$ at shell $m$, the harmonic ladder masses $m_\tau\to m_\mu\to m_e$, the lock-in vev $v=\sqrt{\eta_{\rm paper}\,\Omega_k(m_{\rm lockin};m_{\rm lockin})}$, the proton anchor at $m=\mathrm{referenceM}=4$, the baryon-asymmetry $\eta\propto 6^7\sqrt{3}/(m+1)$, and the rational coefficients $\alpha=3/5$, $\gamma=2/5$, $\alpha_{\rm GUT}=1/42$---taking a literal $m\to\infty$ or sub-Planck continuum limit would \emph{erase} the discrete index that is doing the predictive work.  Continuum forms of these objects are therefore not preferred refinements; they are degenerate, information-losing readouts of the same patch data.

\paragraph{An exception: $\mathfrak{so}(8)$ as a de-facto continuum algebra from a discrete construction.}
Principled exception to the discrete-only stance: the closed Lie algebra $\mathfrak{so}(8) \supseteq G_2\cup\{\Delta\}$ generated from the Fano octonion left-multiplication table and the phase-lift generator $\Delta\in(\mathrm{e}_1,\mathrm{e}_7)$.  Although $\mathfrak{so}(8)$ is a continuous (28-dimensional real) Lie algebra, it is built here entirely from \emph{discrete} ingredients (the seven imaginary octonion units, a finite Fano table, and one $\pi/2$ rotation), and its closure to 28 dimensions is a finite linear-algebra certificate (see the \texttt{Hqiv.SO8Closure} / \texttt{Hqiv.SO8ClosureInterface} stack).  The Lie-group structure is therefore a \emph{de-facto continuum algebra obtained from a discrete construction}: continuous in parameter space, but with no continuum spacetime, no continuum measure, and no UV limit attached.  When this manuscript or its companions write $\exp(t\,\Delta)\in\mathrm{SO}(8)$ or treat $\mathfrak{so}(8)$ rotations as smooth, those are statements internal to a finite-dimensional Lie group whose generators are certified by discrete data; they are \emph{not} a continuum-spacetime commitment and do not require a continuum-QFT scaffolding.

% From papers/paper/*.tex: % Shared reader contract: patch QFT + discrete numerics.
% From papers/*.tex:     % Shared reader contract: patch QFT + discrete numerics.
% From papers/*.tex:     \input{include/patch_theory_messaging.tex}
% From papers/paper/*.tex: \input{../include/patch_theory_messaging.tex}

\paragraph{Patch QFT and discrete numerics (reader contract).}
HQIV (Horizon-Quantized Informational Vacuum) is formulated as a \emph{patch theory}: quantum fields, actions, and physical readouts are attached to discrete patches---shell indices $m\in\mathbb{N}$, finite spacetime index types (e.g.\ $\mathrm{Fin}\,4$), and horizon-limited charts---rather than to a hypothetical fundamental smooth spacetime obtained as a mesh-refinement or ultraviolet limit.
Accordingly, when this text refers to ``QFT,'' it means \emph{patch QFT}: patching rules, finite measurement ledgers, and variational identities on the discrete spine; it is not the claim that the theory is \emph{defined} by recovering ordinary continuum quantum field theory through such a limit.
The numbers that admit the sharpest statements---mass and coupling ladders, curvature ratios, shell thermodynamics, and rational witnesses in \texttt{HQIV\_LEAN}---are objects of \emph{discrete mathematics} on $\mathbb{N}$ (combinatorial growth, cumulative simplex ledgers) together with the ordered field $\mathbb{R}$ as used in the proof assistant for analysis, bounds, and phenomenological display.
Smooth-manifold or continuum field notation, when it appears, is a \emph{translation layer} for comparison with standard physics literature; it does not supersede the discrete definitions.

\paragraph{A continuum is not required, and in places would destroy these results.}
The discrete patch layer is \emph{not} a numerical approximation to a more fundamental continuum theory awaiting future construction; it is the ontology.  Where the literature treats ``the continuum limit'' as the goal, HQIV treats it as a \emph{calculation approximation} for comparison with standard formulas.  The continuum form is an optional convenience for comparison; the proved discrete statement is what carries the physics.  In several load-bearing places---the curvature imprint $\delta_E(m)$ at shell $m$, the harmonic ladder masses $m_\tau\to m_\mu\to m_e$, the lock-in vev $v=\sqrt{\eta_{\rm paper}\,\Omega_k(m_{\rm lockin};m_{\rm lockin})}$, the proton anchor at $m=\mathrm{referenceM}=4$, the baryon-asymmetry $\eta\propto 6^7\sqrt{3}/(m+1)$, and the rational coefficients $\alpha=3/5$, $\gamma=2/5$, $\alpha_{\rm GUT}=1/42$---taking a literal $m\to\infty$ or sub-Planck continuum limit would \emph{erase} the discrete index that is doing the predictive work.  Continuum forms of these objects are therefore not preferred refinements; they are degenerate, information-losing readouts of the same patch data.

\paragraph{An exception: $\mathfrak{so}(8)$ as a de-facto continuum algebra from a discrete construction.}
Principled exception to the discrete-only stance: the closed Lie algebra $\mathfrak{so}(8) \supseteq G_2\cup\{\Delta\}$ generated from the Fano octonion left-multiplication table and the phase-lift generator $\Delta\in(\mathrm{e}_1,\mathrm{e}_7)$.  Although $\mathfrak{so}(8)$ is a continuous (28-dimensional real) Lie algebra, it is built here entirely from \emph{discrete} ingredients (the seven imaginary octonion units, a finite Fano table, and one $\pi/2$ rotation), and its closure to 28 dimensions is a finite linear-algebra certificate (see the \texttt{Hqiv.SO8Closure} / \texttt{Hqiv.SO8ClosureInterface} stack).  The Lie-group structure is therefore a \emph{de-facto continuum algebra obtained from a discrete construction}: continuous in parameter space, but with no continuum spacetime, no continuum measure, and no UV limit attached.  When this manuscript or its companions write $\exp(t\,\Delta)\in\mathrm{SO}(8)$ or treat $\mathfrak{so}(8)$ rotations as smooth, those are statements internal to a finite-dimensional Lie group whose generators are certified by discrete data; they are \emph{not} a continuum-spacetime commitment and do not require a continuum-QFT scaffolding.

% From papers/paper/*.tex: % Shared reader contract: patch QFT + discrete numerics.
% From papers/*.tex:     \input{include/patch_theory_messaging.tex}
% From papers/paper/*.tex: \input{../include/patch_theory_messaging.tex}

\paragraph{Patch QFT and discrete numerics (reader contract).}
HQIV (Horizon-Quantized Informational Vacuum) is formulated as a \emph{patch theory}: quantum fields, actions, and physical readouts are attached to discrete patches---shell indices $m\in\mathbb{N}$, finite spacetime index types (e.g.\ $\mathrm{Fin}\,4$), and horizon-limited charts---rather than to a hypothetical fundamental smooth spacetime obtained as a mesh-refinement or ultraviolet limit.
Accordingly, when this text refers to ``QFT,'' it means \emph{patch QFT}: patching rules, finite measurement ledgers, and variational identities on the discrete spine; it is not the claim that the theory is \emph{defined} by recovering ordinary continuum quantum field theory through such a limit.
The numbers that admit the sharpest statements---mass and coupling ladders, curvature ratios, shell thermodynamics, and rational witnesses in \texttt{HQIV\_LEAN}---are objects of \emph{discrete mathematics} on $\mathbb{N}$ (combinatorial growth, cumulative simplex ledgers) together with the ordered field $\mathbb{R}$ as used in the proof assistant for analysis, bounds, and phenomenological display.
Smooth-manifold or continuum field notation, when it appears, is a \emph{translation layer} for comparison with standard physics literature; it does not supersede the discrete definitions.

\paragraph{A continuum is not required, and in places would destroy these results.}
The discrete patch layer is \emph{not} a numerical approximation to a more fundamental continuum theory awaiting future construction; it is the ontology.  Where the literature treats ``the continuum limit'' as the goal, HQIV treats it as a \emph{calculation approximation} for comparison with standard formulas.  The continuum form is an optional convenience for comparison; the proved discrete statement is what carries the physics.  In several load-bearing places---the curvature imprint $\delta_E(m)$ at shell $m$, the harmonic ladder masses $m_\tau\to m_\mu\to m_e$, the lock-in vev $v=\sqrt{\eta_{\rm paper}\,\Omega_k(m_{\rm lockin};m_{\rm lockin})}$, the proton anchor at $m=\mathrm{referenceM}=4$, the baryon-asymmetry $\eta\propto 6^7\sqrt{3}/(m+1)$, and the rational coefficients $\alpha=3/5$, $\gamma=2/5$, $\alpha_{\rm GUT}=1/42$---taking a literal $m\to\infty$ or sub-Planck continuum limit would \emph{erase} the discrete index that is doing the predictive work.  Continuum forms of these objects are therefore not preferred refinements; they are degenerate, information-losing readouts of the same patch data.

\paragraph{An exception: $\mathfrak{so}(8)$ as a de-facto continuum algebra from a discrete construction.}
Principled exception to the discrete-only stance: the closed Lie algebra $\mathfrak{so}(8) \supseteq G_2\cup\{\Delta\}$ generated from the Fano octonion left-multiplication table and the phase-lift generator $\Delta\in(\mathrm{e}_1,\mathrm{e}_7)$.  Although $\mathfrak{so}(8)$ is a continuous (28-dimensional real) Lie algebra, it is built here entirely from \emph{discrete} ingredients (the seven imaginary octonion units, a finite Fano table, and one $\pi/2$ rotation), and its closure to 28 dimensions is a finite linear-algebra certificate (see the \texttt{Hqiv.SO8Closure} / \texttt{Hqiv.SO8ClosureInterface} stack).  The Lie-group structure is therefore a \emph{de-facto continuum algebra obtained from a discrete construction}: continuous in parameter space, but with no continuum spacetime, no continuum measure, and no UV limit attached.  When this manuscript or its companions write $\exp(t\,\Delta)\in\mathrm{SO}(8)$ or treat $\mathfrak{so}(8)$ rotations as smooth, those are statements internal to a finite-dimensional Lie group whose generators are certified by discrete data; they are \emph{not} a continuum-spacetime commitment and do not require a continuum-QFT scaffolding.


\paragraph{Patch QFT and discrete numerics (reader contract).}
HQIV (Horizon-Quantized Informational Vacuum) is formulated as a \emph{patch theory}: quantum fields, actions, and physical readouts are attached to discrete patches---shell indices $m\in\mathbb{N}$, finite spacetime index types (e.g.\ $\mathrm{Fin}\,4$), and horizon-limited charts---rather than to a hypothetical fundamental smooth spacetime obtained as a mesh-refinement or ultraviolet limit.
Accordingly, when this text refers to ``QFT,'' it means \emph{patch QFT}: patching rules, finite measurement ledgers, and variational identities on the discrete spine; it is not the claim that the theory is \emph{defined} by recovering ordinary continuum quantum field theory through such a limit.
The numbers that admit the sharpest statements---mass and coupling ladders, curvature ratios, shell thermodynamics, and rational witnesses in \texttt{HQIV\_LEAN}---are objects of \emph{discrete mathematics} on $\mathbb{N}$ (combinatorial growth, cumulative simplex ledgers) together with the ordered field $\mathbb{R}$ as used in the proof assistant for analysis, bounds, and phenomenological display.
Smooth-manifold or continuum field notation, when it appears, is a \emph{translation layer} for comparison with standard physics literature; it does not supersede the discrete definitions.

\paragraph{A continuum is not required, and in places would destroy these results.}
The discrete patch layer is \emph{not} a numerical approximation to a more fundamental continuum theory awaiting future construction; it is the ontology.  Where the literature treats ``the continuum limit'' as the goal, HQIV treats it as a \emph{calculation approximation} for comparison with standard formulas.  The continuum form is an optional convenience for comparison; the proved discrete statement is what carries the physics.  In several load-bearing places---the curvature imprint $\delta_E(m)$ at shell $m$, the harmonic ladder masses $m_\tau\to m_\mu\to m_e$, the lock-in vev $v=\sqrt{\eta_{\rm paper}\,\Omega_k(m_{\rm lockin};m_{\rm lockin})}$, the proton anchor at $m=\mathrm{referenceM}=4$, the baryon-asymmetry $\eta\propto 6^7\sqrt{3}/(m+1)$, and the rational coefficients $\alpha=3/5$, $\gamma=2/5$, $\alpha_{\rm GUT}=1/42$---taking a literal $m\to\infty$ or sub-Planck continuum limit would \emph{erase} the discrete index that is doing the predictive work.  Continuum forms of these objects are therefore not preferred refinements; they are degenerate, information-losing readouts of the same patch data.

\paragraph{An exception: $\mathfrak{so}(8)$ as a de-facto continuum algebra from a discrete construction.}
Principled exception to the discrete-only stance: the closed Lie algebra $\mathfrak{so}(8) \supseteq G_2\cup\{\Delta\}$ generated from the Fano octonion left-multiplication table and the phase-lift generator $\Delta\in(\mathrm{e}_1,\mathrm{e}_7)$.  Although $\mathfrak{so}(8)$ is a continuous (28-dimensional real) Lie algebra, it is built here entirely from \emph{discrete} ingredients (the seven imaginary octonion units, a finite Fano table, and one $\pi/2$ rotation), and its closure to 28 dimensions is a finite linear-algebra certificate (see the \texttt{Hqiv.SO8Closure} / \texttt{Hqiv.SO8ClosureInterface} stack).  The Lie-group structure is therefore a \emph{de-facto continuum algebra obtained from a discrete construction}: continuous in parameter space, but with no continuum spacetime, no continuum measure, and no UV limit attached.  When this manuscript or its companions write $\exp(t\,\Delta)\in\mathrm{SO}(8)$ or treat $\mathfrak{so}(8)$ rotations as smooth, those are statements internal to a finite-dimensional Lie group whose generators are certified by discrete data; they are \emph{not} a continuum-spacetime commitment and do not require a continuum-QFT scaffolding.


\paragraph{Patch QFT and discrete numerics (reader contract).}
HQIV (Horizon-Quantized Informational Vacuum) is formulated as a \emph{patch theory}: quantum fields, actions, and physical readouts are attached to discrete patches---shell indices $m\in\mathbb{N}$, finite spacetime index types (e.g.\ $\mathrm{Fin}\,4$), and horizon-limited charts---rather than to a hypothetical fundamental smooth spacetime obtained as a mesh-refinement or ultraviolet limit.
Accordingly, when this text refers to ``QFT,'' it means \emph{patch QFT}: patching rules, finite measurement ledgers, and variational identities on the discrete spine; it is not the claim that the theory is \emph{defined} by recovering ordinary continuum quantum field theory through such a limit.
The numbers that admit the sharpest statements---mass and coupling ladders, curvature ratios, shell thermodynamics, and rational witnesses in \texttt{HQIV\_LEAN}---are objects of \emph{discrete mathematics} on $\mathbb{N}$ (combinatorial growth, cumulative simplex ledgers) together with the ordered field $\mathbb{R}$ as used in the proof assistant for analysis, bounds, and phenomenological display.
Smooth-manifold or continuum field notation, when it appears, is a \emph{translation layer} for comparison with standard physics literature; it does not supersede the discrete definitions.

\paragraph{A continuum is not required, and in places would destroy these results.}
The discrete patch layer is \emph{not} a numerical approximation to a more fundamental continuum theory awaiting future construction; it is the ontology.  Where the literature treats ``the continuum limit'' as the goal, HQIV treats it as a \emph{calculation approximation} for comparison with standard formulas.  The continuum form is an optional convenience for comparison; the proved discrete statement is what carries the physics.  In several load-bearing places---the curvature imprint $\delta_E(m)$ at shell $m$, the harmonic ladder masses $m_\tau\to m_\mu\to m_e$, the lock-in vev $v=\sqrt{\eta_{\rm paper}\,\Omega_k(m_{\rm lockin};m_{\rm lockin})}$, the proton anchor at $m=\mathrm{referenceM}=4$, the baryon-asymmetry $\eta\propto 6^7\sqrt{3}/(m+1)$, and the rational coefficients $\alpha=3/5$, $\gamma=2/5$, $\alpha_{\rm GUT}=1/42$---taking a literal $m\to\infty$ or sub-Planck continuum limit would \emph{erase} the discrete index that is doing the predictive work.  Continuum forms of these objects are therefore not preferred refinements; they are degenerate, information-losing readouts of the same patch data.

\paragraph{An exception: $\mathfrak{so}(8)$ as a de-facto continuum algebra from a discrete construction.}
Principled exception to the discrete-only stance: the closed Lie algebra $\mathfrak{so}(8) \supseteq G_2\cup\{\Delta\}$ generated from the Fano octonion left-multiplication table and the phase-lift generator $\Delta\in(\mathrm{e}_1,\mathrm{e}_7)$.  Although $\mathfrak{so}(8)$ is a continuous (28-dimensional real) Lie algebra, it is built here entirely from \emph{discrete} ingredients (the seven imaginary octonion units, a finite Fano table, and one $\pi/2$ rotation), and its closure to 28 dimensions is a finite linear-algebra certificate (see the \texttt{Hqiv.SO8Closure} / \texttt{Hqiv.SO8ClosureInterface} stack).  The Lie-group structure is therefore a \emph{de-facto continuum algebra obtained from a discrete construction}: continuous in parameter space, but with no continuum spacetime, no continuum measure, and no UV limit attached.  When this manuscript or its companions write $\exp(t\,\Delta)\in\mathrm{SO}(8)$ or treat $\mathfrak{so}(8)$ rotations as smooth, those are statements internal to a finite-dimensional Lie group whose generators are certified by discrete data; they are \emph{not} a continuum-spacetime commitment and do not require a continuum-QFT scaffolding.


%============================================================================
\section{Scope and claim discipline}\label{sec:scope}

This paper does not reprove the Standard Model. It reads the
electroweak scale and a fermion mass ladder off horizon
quantities already proved upstream, in the discrete-patch
ontology defined by the patch-theory reader contract above.

\begin{enumerate}[leftmargin=2em]
\item \textbf{Two upstream axioms.} The discrete null lattice
  with Planck cell size and informational monogamy
  ($\alpha=3/5$, hence $\gamma=2/5$) are imported as black-box
  certificates from the lightcone paper
  \cite{hqiv-lightcone-paper} and the closure paper
  \cite{hqiv-so8-paper}. The
  $G_2\cup\{\Delta\}\Rightarrow\mathfrak{so}(8)$ Lie closure
  used as bookkeeping for representation labels is the same
  certificate.
\item \textbf{One upstream variational layer.} The O-Maxwell
  action, its discrete Euler--Lagrange identity, the
  cyclic-Stokes identity for plaquette transports, and the
  HQVM gravitational constraint action are imported as
  black-box theorems from the octonionic-action paper
  \cite{hqiv-oct-action-paper}. The weak/Higgs scaffold of
  Section~\ref{sec:core}(III) is layered on top of that, with
  proved collapse back to the abelian core.
\item \textbf{One upstream finite-mode spectrum.} The
  finite-mode Kirchhoff law, the proton lock-in scale-witness,
  and the wall-clock vs apparent-age discriminator are
  imported as black-box theorems from the finite-mode
  Kirchhoff paper \cite{hqiv-kirchhoff-paper}. This paper
  reuses the same single-witness discipline.
\item \textbf{One active scale witness per solve.} Throughout
  this paper the active witness is the
  \hyperref[lean:proton-lockin]{proton lock-in scale-witness},
  the chart-pinning convention also used in
  \cite{hqiv-kirchhoff-paper}. CODATA $1/\alpha$
  \cite{codata2018}, CMB horizons, and PDG quark/lepton
  centrals \cite{pdg2024} are predictions or comparison
  layers, not simultaneous anchors.
\item \textbf{No continuum-limit programme.} Per the
  patch-theory reader contract, the discrete patch layer is
  the ontology. The $\mathbb{C}$-valued
  $\mathrm{Fin}\,8$ carrier and the $\mathfrak{so}(8)$
  Lie-algebra layer are de-facto continuum objects built from
  finite, discrete data; there is no continuum-spacetime
  limit, no continuum-QFT recovery, and no mesh-refinement
  programme anywhere in this chain.
\item \textbf{Honesty about quark anchors.} The quark ladder
  in Section~\ref{sec:quarks} still carries explicit GeV
  anchors and $\mathbb{N}$ shell-coordinate triples. These
  are flagged in-line; what is and is not derived from
  $(T_{\mathrm{lockin}},S,\gamma,\alpha,M)$ is itemised
  there. The proton mass under the lock-in witness, by
  contrast, is the chart-pinning anchor: a non-trivial
  HQIV calculation at $M=4$, but not a two-axiom output.
\end{enumerate}

%============================================================================
\section{The core derivation}
\label{sec:core}

\paragraph{(I) Electroweak scale from $T_{\mathrm{lockin}}$ and outer horizons.}
Let $m_{\mathrm{lockin}}=M=\mathrm{referenceM}$ and
$T_{\mathrm{lockin}}=T(M)=1/(M+1)$ on the
\hyperref[lean:auxiliary-T]{auxiliary-field temperature ladder}.
The
\hyperref[lean:gauge-lepton]{derived gauge-and-lepton sector
module} defines a geometric vacuum scale
\[
  v \;=\; T_{\mathrm{lockin}}\cdot S(M+1)\cdot(1+\gamma),
  \qquad S(m)=(m+1)(m+2),
\]
with $\gamma=1-\alpha=\tfrac{2}{5}$. Gauge and scalar masses are
built from $v$ with the
\hyperref[lean:gauge-couplings]{triality-fixed couplings}
($\mathrm{SU}(2)_w$ and $\mathrm{U}(1)_Y$ derived couplings).
The module proves closed rational witnesses---for example
$M_W^{\mathrm{derived}}=392/5$ in the module's natural units, with
matching $M_Z$ and $m_H$ identities (the
\hyperref[lean:boson-witnesses]{rational-boson witness pack})---
together with comparison theorems against electroweak reference
values \cite{pdg2024}. Neutrino masses begin from the
\hyperref[lean:neutrino-ladder]{outer-horizon neutrino mass
identity}
\(
  m_{\nu_e}^{\mathrm{derived}}=\bigl(\gamma/S(M+2)\bigr)\,M_Z^{\mathrm{derived}}
\)
with a generation cascade to the higher families. \emph{No PDG
literals are imported to define these objects.}

\paragraph{(II) Combinatorial imprint and horizon-surface readouts.}
The upstream lightcone paper \cite{hqiv-lightcone-paper} proves
the
\hyperref[lean:alpha-ratio]{lattice $\alpha$-ratio identity}
$\alpha=3/5$ at every truncation. The
\hyperref[lean:fano-resonance]{Fano resonance package} provides
the shared surface $S(m)$ and the
\hyperref[lean:detuned-shell]{detuned outer-shell surface}; mass
ratios are ratios of these, giving an $S(m)$-octave hierarchy that
feeds both the
\hyperref[lean:lepton-resonance]{charged-lepton resonance ladder}
and the
\hyperref[lean:quark-meta]{quark meta-resonance ladder}. Where a
content-class ansatz is needed, the
\hyperref[lean:content-mass-bridge]{conserved-content
mass-scaling bridge} formalises the expected rank-$l$ scaling on
top of the shell imprint---the combinatorial hierarchy is
first-class, not an afterthought.

The preferred lepton path now treats the modal frequency /
horizon data as primary:
\[
  \text{(modal-frequency horizon)} \;\longrightarrow\;
  \text{(global-detuning lepton ratios)} \;\longrightarrow\;
  \text{(derived gauge-and-lepton sector)},
\]
with the
\hyperref[lean:modal-horizon]{modal-frequency horizon module},
\hyperref[lean:lepton-global-detuning]{global-detuning lepton
module}, and
\hyperref[lean:gauge-lepton]{derived gauge-and-lepton sector
module} respectively (and the
\hyperref[lean:lepton-resonance]{charged-lepton resonance
module} retained as a thicker phenomenology layer). This
reflects the project stance that horizons and effective
surfaces are primary, while shell numbers are sampled readouts.

\paragraph{(III) Variational weak/Higgs scaffold on the same lock-in channel.}
Beyond the mass-ladder exports, the
\hyperref[lean:weak-higgs]{weak/Higgs O-Maxwell scaffold}
provides a typed extension interface on the proved O-Maxwell
action of the octonionic-action paper
\cite{hqiv-oct-action-paper}:
\[
  F^{a,\mathrm{weak}}_{\mu\nu}
  =
  (W^a_\nu-W^a_\mu)+g_w\,\mathrm{comm}^a_{\mu\nu},
  \qquad
  \mathcal{L}_{\mathrm{weak}}
  =-\frac14\sum_{a,\mu,\nu}\frac{(F^{a,\mathrm{weak}}_{\mu\nu})^2}{2},
\]
with scalar potential $\lambda(|\Phi|^2-v^2)^2$ on
$\Phi\in\mathbb{R}^8$. The lock-in vev proxy is the
\hyperref[lean:lockin-vev]{lock-in vev identity}
\[
  v_{\mathrm{lockin}}=\sqrt{\eta_{\mathrm{paper}}\,\Omega_k(M;M)}.
\]
The variational scaffold proves reduction and normalisation
identities that go beyond placeholder content (collected in the
\hyperref[lean:weak-higgs-reductions]{weak/Higgs reduction lemma
pack}): the weak field strength reduces to its abelian form on
the trivial commutator, the diagonal of $F^{a,\mathrm{weak}}$
vanishes there, the lock-in vev squares to
$\eta_{\mathrm{paper}}$, and the combined weak+Higgs action
reduces exactly to the abelian O-Maxwell action under those
specialisations. Status: abelian action dynamics are proved in
the upstream octonionic-action paper
\cite{hqiv-oct-action-paper}; the weak/Higgs layer is a formally
controlled extension surface with proved collapse back to the
abelian core and proved lock-in / triality consistency
identities.

%============================================================================
\section{Maxwell spectrum and the \texorpdfstring{$\Delta$}{Delta}-triality spine}
\label{sec:maxwell-delta-triality}

The detuning denominator used in the active mass-ladder path is
not a free parameter; it is sourced from a matrix-level
O-Maxwell / Fano construction (the
\hyperref[lean:fano-spectrum]{Fano O-Maxwell spectrum module}).
The construction is:

\begin{enumerate}[leftmargin=2em]
  \item \textbf{Carrier and line projector.} A Fano vertex is
  embedded into the octonion matrix carrier $\mathrm{Fin}\,8$
  as index $1..7$ (scalar slot $0$ reserved), then lifted to a
  diagonal selector. Summing selectors over a line $L$ gives
  \[
    \Pi_L \;=\; \sum_{v\in L}\Pi_v,
    \qquad \operatorname{tr}(\Pi_L)=3
  \]
  (the
  \hyperref[lean:fano-line-trace]{Fano-line selector-trace
  identity}).

  \item \textbf{Phase-lift generator $\Delta$.} The
  \hyperref[lean:phase-lift]{phase-lift generator module} fixes
  the antisymmetric $8\times 8$ generator with nonzero
  $(1,7)$, $(7,1)$ entries; the scaffold proves
  \[
    \operatorname{tr}(\Delta\Delta^{\mathsf T})=2
  \]
  (the
  \hyperref[lean:phase-lift-quadratic]{phase-lift quadratic
  trace identity}).

  \item \textbf{Projection normalisation.} The spectral scalar
  is defined by
  \[
    \mathcal{N}(L)=\frac{\operatorname{tr}(\Pi_L)+\operatorname{tr}(\Delta\Delta^{\mathsf T})}{5}
    =\frac{3+2}{5}=1
  \]
  (the
  \hyperref[lean:fano-spectral-norm]{spectral projection
  normalisation theorem}).

  \item \textbf{Projected mode strength and first-order jet.}
  For mode $(L,m)$, the projected strength is $\mathcal{N}(L)$
  times the
  \hyperref[lean:phase-lift-coeff]{phase-lift coefficient}
  $\varphi(m)/6$, with $\varphi(m)=2(m+1)$. Combining
  $\mathcal{N}(L)=1$ with the standard Rindler-detuning
  source then gives the
  \hyperref[lean:fano-rindler-1jet]{first-order Fano--Rindler
  jet identity}
  \[
    D_L(m)=1+\frac{\gamma}{2}\,m,
  \]
  i.e.\ exactly the
  \hyperref[lean:rindler-shared]{shared Rindler detuning}
  used downstream.
\end{enumerate}

The public detuning denominator (the
\hyperref[lean:triality-denominator]{triality-projected
denominator}) coincides with this spectral 1-jet, and the
companion
\hyperref[lean:triality-denominator-firstorder]{first-order
identity} is therefore a proved corollary, not a placeholder.

\paragraph{What ``triality'' means here today.}
The
\hyperref[lean:triality-cycle]{Spin(8) triality module} proves
the order-3 representation cycle on
$(8_v,8_{s^+},8_{s^-})$. What is \emph{not} yet proved at this
revision is the stronger forcing theorem that the denominator's
unit constant is uniquely imposed by triality-equivariant mode
selection; that statement is itemised under open theorem
objectives in Section~\ref{sec:roadmap}.

%============================================================================
\section{Backbone definitions (one pass)}
\label{sec:backbone}

The backbone of the construction is a small handful of horizon
quantities:
\begin{itemize}[leftmargin=2em]
\item the
  \hyperref[lean:auxiliary-T]{auxiliary-field temperature
  ladder} $T(m)=1/(m+1)$, the harmonic-ladder content
  $\varphi(m)=2(m+1)$, and the shell-shape constants;
\item the
  \hyperref[lean:phase-lift]{phase-lift generator} (fixed
  antisymmetric $\Delta$) and the phase-lift coefficient
  $\varphi(m)/6$;
\item the
  \hyperref[lean:triality-cycle]{Spin(8) triality module} for
  representation-cycle bookkeeping;
\item the
  \hyperref[lean:sm-gr]{SM--GR unification module} which
  fixes $\alpha_{\mathrm{GUT}}=1/42$ from the cube directions
  $\times$ octonion imaginary units;
\item the
  \hyperref[lean:baryogenesis]{baryogenesis module} (lock-in
  index $m_{\mathrm{lockin}}=M$); and
\item the
  \hyperref[lean:bound-states]{bound-states module} together
  with the
  \hyperref[lean:harmonic-ladder]{harmonic-ladder mass
  module} for shell-resolved $\alpha_{\mathrm{eff}}(m)$ and
  hydrogenic $\alpha^2$ scalings.
\end{itemize}

For the variational extension path, two further modules enter:
the upstream
\hyperref[lean:action]{O-Maxwell action module} (whose
abelian Euler--Lagrange identity is the proved base
\cite{hqiv-oct-action-paper}) and the
\hyperref[lean:weak-higgs]{weak/Higgs O-Maxwell scaffold}, with
its
\hyperref[lean:weak-higgs-reductions]{reduction-lemma pack}
proving collapse back to the abelian core.

%============================================================================
\section{Fermion and hadron exports}
\label{sec:fermion}

The
\hyperref[lean:sm-gr]{SM--GR unification module} exports
Planck-unit masses via the
\hyperref[lean:sm-mass]{geometric SM mass functional}: a
$\tau$-Planck scale divided by the
\hyperref[lean:resonance-product]{resonance-product factor}
defined in the
\hyperref[lean:lepton-resonance]{charged-lepton resonance
module}. The heavy triality slot has resonance product equal to
$1$, so the $\tau$, top, and bottom labels coincide at that
functional's value (a structural feature of the triality
quotient, not a coincidence in the data).

%============================================================================
\section{Single-scale witness and informational-energy mass row}
\label{sec:scale-witness}

\paragraph{One active witness per pipeline.}
The
\hyperref[lean:scale-witness]{scale-witness enum} carries an
inductive type with three modes: the
\hyperref[lean:proton-lockin]{proton lock-in scale-witness}
(default), the CODATA $\alpha$ witness, and the present-CMB
witness. The discipline is operational, not a new Mathlib
axiom: in any given export or Python solve, \textbf{exactly
one} of these may pin the dimensionful chart; the others are
cross-checks.
\begin{itemize}[leftmargin=2em]
  \item \textbf{Proton lock-in (default).} The
    \hyperref[lean:derived-proton]{derived proton mass} at
    $M=\mathrm{referenceM}$ fixes the mass/unit chart; the
    neutron mass follows from the same lock-in shell and
    constituent content (the
    \hyperref[lean:delta-m]{neutron--proton gap identity},
    udd vs.\ uud); $W/Z/H$ masses load from the Lean closure
    witnesses in the
    \hyperref[lean:witnesses-json]{HQIV witness JSON bundle}.
    CODATA $1/\alpha$ \cite{codata2018} is a \emph{predicted}
    EM coupling, not a row in the solve.
  \item \textbf{CODATA $\alpha$ witness.} A regression mode
    used during cross-checks: the continuous Gauss$\to$EW brace
    pins $1/\alpha$ to CODATA while the rest of the geometry
    supplies the remaining rows.
  \item \textbf{Present-CMB witness.} A shallow-$\xi$
    comparison chart ($\xi\approx 1.07$ mean-field), not the
    brace mass row at lock-in ($\xi=5$).
\end{itemize}
The
\hyperref[lean:hqiv-witnesses]{HQIV witnesses Lake target}
exports the JSON bundle that downstream Python loaders read.

\paragraph{Informational energy and readout gauges.}
In natural units ($c=\hbar=T_{\mathrm{Pl}}=1$), the
\hyperref[lean:info-energy]{informational-energy module}
defines
\[
  E_{\mathrm{tot}}(m_{\mathrm{rest}},\xi) = m_{\mathrm{rest}} + \frac{1}{\Theta_{\mathrm{local}}(\xi)},
  \qquad \Theta_{\mathrm{local}}(\xi)=\frac{T_{\mathrm{Pl}}}{\xi},
  \qquad \mathrm{loc}(\xi)=\frac{1}{\Theta_{\mathrm{local}}(\xi)}=\xi.
\]
Bosons use the
\hyperref[lean:additive-loc]{additive-localisation readout
gauge} (full $E_{\mathrm{tot}}$ is the observable mass);
hadrons use the
\hyperref[lean:lapse-readout]{multiplicative-lapse readout
gauge} (rest slot $\div N_{\mathrm{lapse}}$ with localisation
carried by the lapse increment, not double-counted). The
\hyperref[lean:gauge-transform]{gauge-transformation lemma}
records when the two readout gauges agree after calibrating
$m_{\mathrm{rest}}$.

\paragraph{Mass row on the Fano coupling system.}
The
\hyperref[lean:cont-xi-coupling]{continuous-$\xi$ coupling
module} packages the Fano coupling system rows in Lean form.
The \emph{default} mass-sector row is
\[
  c_0 + \mathrm{loc}(\xi_G) = 2\pi\,\Omega_k(\xi_G),
  \qquad
  \Omega_k(\xi)=\frac{I(\xi)}{I(\xi_{\mathrm{lock}})},
\]
with $\xi_{\mathrm{lock}}=M+1=5$ so
$\Omega_k(\xi_{\mathrm{lock}})=1$. The
\hyperref[lean:info-energy-row]{informational-energy mass-row
identity} packages the budget identity; when the row holds, the
\hyperref[lean:info-energy-row-twopi]{$2\pi\,\Omega_k$
specialisation} gives $E_{\mathrm{tot}}=2\pi\,\Omega_k(\xi_G)$.

\paragraph{Three $\Omega_k$ charts (do not conflate).}
Three distinct $\Omega_k$ readouts must be kept separate:
\begin{itemize}[leftmargin=2em]
\item the \emph{brace/coupling} value $\Omega_k(\xi_G)\approx 0.72$ at
  $\xi_G\approx 3.47$ with $\xi_{\mathrm{lock}}=5$;
\item the \emph{shallow} $\Omega_k(\xi\approx 1.07)\approx 0.03$
  (not the brace root);
\item the \emph{$\Omega_k^{\mathrm{true}}$} axiom limit
  $\approx 0.0098$ and CMB-as-horizon scales $\sim 10^{-28}$---
  comparison layers only, not slots in the brace mass row.
\end{itemize}

\paragraph{Excited states.}
The
\hyperref[lean:meta-horizon]{meta-horizon excited-states
module} models baryon excitations as radial ($n$) and orbital
($\ell$) modes on the lock-in drum:
\(
  \mathrm{totalModeMass}(n,\ell) =
  M_{\mathrm{constituent}}-E_{\mathrm{bind}}(M+n+\ell)
\)
with the same nucleon composite-trace witness as the ground
state. The downstream calculator uses a ``ground + $\Delta M$''
breakdown: an informational-energy ground mass (no excitation
tag), plus $\Delta M_{\mathrm{radial}}$ for decuplet tags and
$\Delta M_{\mathrm{orbital}}$ for vector tags. \emph{Open:} the
raw composite-trace binding grows with shell index (wrong sign
for excitation energy) at the current revision; the operational
$\Delta M$ uses the meta-horizon surface step
$m_p\cdot(S_{m+1}/S_m-1)$ until the binding shell law is
refined and exported.

%============================================================================
\section{Quarks and nucleons: what is derived, what is still a witness}
\label{sec:quarks}

\paragraph{Honesty (not ``no witnesses'').}
The
\hyperref[lean:quark-meta]{quark meta-resonance ladder} is
\emph{not} free of external numerals. In particular:
\begin{itemize}[leftmargin=2em]
  \item \textbf{Absolute scale.} The up-type ladder is
    normalised by a named real \texttt{m\_top\_GeV} (a
    PDG-style pole-mass listing \cite{pdg2024} used as a
    calibration witness). Ratios between generations are
    \emph{not} taken from PDG tables; they are fixed by
    \hyperref[lean:resonance-step]{geometric-resonance-step
    ratios} built from
    \hyperref[lean:detuned-shell]{detuned outer-shell
    surfaces} and explicit $\mathbb{N}$ shell endpoints.
  \item \textbf{Readout tables.} The integer triples used as
    \texttt{m\_quark\_up\_*} and \texttt{m\_quark\_down\_*}
    shell endpoints are calibration inputs that determine
    internal resonance factors; uniqueness of these
    coordinates from the null lattice alone is \emph{not}
    proved in this module.
  \item \textbf{Legacy down branch.} A separate real
    \texttt{m\_bottom\_GeV} still fixes the old down-type
    comparison ladder unless and until that path is replaced
    end-to-end by the active API.
\end{itemize}

\paragraph{Lock-in readout vs.\ GeV top anchor (disambiguation).}
The lock-in indices agree across modules: in the
\hyperref[lean:baryogenesis]{baryogenesis module} and the
\hyperref[lean:quark-meta]{quark meta-resonance module},
$m_{\mathrm{lockin}}=M=\mathrm{referenceM}$ is also the index
$m_{\mathrm{top\,at\,lockin}}$, with
\(
  T_{\mathrm{lockin}} = T(m_{\mathrm{top\,at\,lockin}})
\)
on the
\hyperref[lean:auxiliary-T]{auxiliary-field temperature ladder}
(so $T(M)=1/(M+1)$ in natural units, not a GeV number). The
\hyperref[lean:t-lockin-id]{lock-in temperature identity}
records this. So the story is: lock-in, baryogenesis, and the
``top-birth'' readout are the \emph{same} calibration row, and
$T_{\mathrm{lockin}}$ is the horizon temperature assigned to
that row. That is not the same proposition as ``the real number
$T_{\mathrm{lockin}}$ equals the pole mass $m_t$ in GeV'': at
$M=4$ one has $T=1/5$ in those units, not $\sim 172\,\mathrm{GeV}$.
The explicit real $\mathtt{m\_top\_GeV}$ exists precisely to
supply \textbf{laboratory GeV} normalisation for the
\hyperref[lean:color-resonance]{color-resonance mass ladder};
it is the witness that remains separate from the abstract
ladder unless a future theorem closes the gap. The
\hyperref[lean:sm-mass]{geometric SM mass functional} carries a
parallel export $m_t^{\mathrm{natural}}$ that uses the $\tau$
Planck mass and resonance products only---no import of
$\mathtt{m\_top\_GeV}$ there. ADM lapse and plasma corrections
\cite{ettinger2026hqiv} are \emph{not} applied inside $T$
itself, so ``top energy'' in the physical paper sense may
require that layer; the Lean lock-in identity is the lock-in
temperature at that readout.

\paragraph{Why top is the heavy anchor (current implementation).}
The active color-composed API fixes the heavy up-like channel
to the top anchor at the lock-in shell. This is not arbitrary
bookkeeping: it matches the lock-in-heavy channel where the
resonance product is $1$, so the heavy band is the unique
non-relaxed reference from which mid/light bands are generated
by the same detuned-surface readout map. In other words, top is
used as the anchor because the current ladder is implemented as
\emph{``heavy lock-in witness $+$ horizon/surface relaxations''}
rather than as three independently normalised flavour masses.
What \emph{is} derived in-structure from the same spine as
leptons ($S(m)$, Rindler detuning, $\gamma$, $\alpha$, $M$) are
the ratio factors between shells and the composite-trace
binding (the
\hyperref[lean:shared-binding]{shared-binding nucleon energy})
at $M$.

\paragraph{Active visible-state API (colour-composed ladder).}
The public masses (the
\hyperref[lean:color-resonance]{color-resonance mass ladder})
normalise the heavy up-like band to the top anchor, then
compress the heavy down-like visible mass by a cross-channel
detuning ratio of detuned outer-shell surfaces, plus a visible
bookkeeping factor $2\times 3$ (two visible charge signs across
three colour channels). Mid and light bands descend by the same
up- and down-resonance-product machinery as elsewhere. Nucleon
constituents use the active light up- and down-quark masses
with a shared dressing scale and a smaller residual detuning
tied to $\alpha_{\mathrm{eff}}$ at $M$; the
\hyperref[lean:derived-nucleons]{derived nucleon-mass module}
exports the proton mass, the neutron mass, and the
neutron--proton gap. Under the proton lock-in scale-witness
(Section~\ref{sec:scale-witness}), the exported proton mass is
the \textbf{primary} mass witness; the
\hyperref[lean:proton-anchor]{proton-anchor mass constant}
\(M_p^{\mathrm{anchor}}=938.272\,\mathrm{MeV}\) remains a
legacy reference row for the
\hyperref[lean:lapse-readout]{multiplicative-lapse readout
gauge}, not a second simultaneous anchor in the downstream
pipeline. The neutron--proton gap is content-only at shared
lock-in $\xi$: udd vs.\ uud constituent energies with the
shared-binding piece cancelling in $\Delta m$.

\paragraph{Representative numerics (module units, current witness set).}
Evaluating the definitions with the literals in the repository
(same formulas as the
\hyperref[lean:quark-meta]{quark meta-resonance module};
approximate) gives:
\begin{center}
\begin{tabular}{l r}
  charm (from top / first up-step) & $\approx 1.27\,\mathrm{GeV}$ \\
  active up-like light             & $\approx 0.00220\,\mathrm{GeV}$ \\
  active down-like heavy (compressed) & $\approx 4.88\,\mathrm{GeV}$ \\
  active down-like light           & $\approx 0.00550\,\mathrm{GeV}$ \\
  derived proton mass              & $\approx 938.3\,\mathrm{MeV}$ \\
  derived neutron mass             & $\approx 939.6\,\mathrm{MeV}$ \\
  neutron--proton gap              & $\approx 1.29\,\mathrm{MeV}$
\end{tabular}
\end{center}
These are \textbf{not} PDG matches by construction; they are
the output of the closed Lean definitions given the current
anchors and readout tables.

\paragraph{Why they differ from observation.}
PDG values \cite{pdg2024} for quarks and hadrons are quoted in
different schemes (pole vs.\ $\overline{\mathrm{MS}}$, running
scales) and include nonperturbative QCD and finite-volume /
lattice input where applicable. This formalisation instead
implements a \textbf{finite-chart resonance surrogate}: fixed
$\mathbb{N}$ readout endpoints, a single heavy-quark GeV
anchor, and a composite-trace binding ansatz at $M$. Differences
are therefore expected:
\begin{itemize}[leftmargin=2em]
  \item \textbf{Proton $\sim 938\,\mathrm{MeV}$.} The current
    bottom-up network export lands near the laboratory anchor
    when the constituent/binding split and active
    light-quark readouts are applied at $M$; this is still a
    finite-chart surrogate, not a lattice-QCD calculation.
    Residual scheme dependence (pole vs.\ $\overline{\mathrm{MS}}$,
    running scale) remains outside the Lean artefact.
  \item \textbf{Light-quark ``masses''.} Experimental
    light-quark masses are indirect and scheme-dependent; our
    GeV numbers are internal ladder coordinates, not a claim
    of pole masses for $u,d$.
  \item \textbf{Heavy flavours.} Charm / bottom / top
    comparisons would require either removing the top anchor
    or deriving it from a deeper lock-in theorem; until then,
    agreement with PDG \cite{pdg2024} is not the target of
    the proof artefact.
\end{itemize}

%============================================================================
\section{Explicit inputs and proof obligations}
\label{sec:roadmap}

The HQIV-LEAN repository's
\texttt{AGENTS/MASS\_DERIVATION\_ROADMAP.md} is the
authoritative status board. In short:

\begin{itemize}[leftmargin=2em]
  \item \textbf{Faithful to the lightcone spine today.}
    Bosonic closure witnesses and the neutrino ladder in the
    \hyperref[lean:gauge-lepton]{derived gauge-and-lepton
    sector module}; structural binding quantities (the
    \hyperref[lean:shared-binding]{shared-binding nucleon
    energy} and composite-trace packaging) tied to $M$,
    $\alpha$, $\gamma$, and $S(m)$; informational-energy
    readout gauges and the $2\pi\,\Omega_k$ mass row in the
    \hyperref[lean:info-energy]{informational-energy module}
    and the
    \hyperref[lean:cont-xi-coupling]{continuous-$\xi$
    coupling module} (Section~\ref{sec:scale-witness}).
  \item \textbf{Single-scale export discipline.} The
    \hyperref[lean:scale-witness]{scale-witness enum} and the
    \hyperref[lean:hqiv-witnesses]{HQIV witnesses Lake
    target} together with the
    \hyperref[lean:witnesses-json]{HQIV witness JSON bundle};
    default active witness is the
    \hyperref[lean:proton-lockin]{proton lock-in
    scale-witness}, and CODATA $\alpha$ \cite{codata2018} and
    PDG masses \cite{pdg2024} are comparison targets unless
    the CODATA-witness mode is explicitly selected.
  \item \textbf{Variational extension formalised as scaffold.}
    The
    \hyperref[lean:weak-higgs]{weak/Higgs O-Maxwell scaffold}
    defines a non-abelian-form weak field strength, Higgs
    potential, lock-in vev proxy, and combined action slot,
    with proved reduction identities back to the abelian
    O-Maxwell action of \cite{hqiv-oct-action-paper}.
  \item \textbf{Still carried by explicit numerals or
    boundary conditions.} Substrate pins (the
    \hyperref[lean:qcd-shell]{QCD-onset shell pin}); the
    quark ladder's GeV anchors (\texttt{m\_top\_GeV}, legacy
    \texttt{m\_bottom\_GeV}) and $\mathbb{N}$ shell triples
    in the
    \hyperref[lean:quark-meta]{quark meta-resonance module}
    (Section~\ref{sec:quarks}); several Standard-Model
    alignment decimals in the
    \hyperref[lean:sm-gr]{SM--GR unification module} where
    the code matches paper fiducials by definition.
  \item \textbf{What the roadmap asks next.}
    Reduce the quark anchor count and unify the legacy
    down-ladder with the active visible-state API; derive or
    eliminate \texttt{m\_top\_GeV} without breaking the
    colour-composed API; on the variational side, promote
    the weak/Higgs scaffold from reduction-level theorems to
    full Euler--Lagrange closure on the weak and scalar
    channels.
\end{itemize}

\noindent This division is intentional: the formalisation
\emph{separates} closure witnesses derived from
$(T_{\mathrm{lockin}},S,\gamma,\alpha,M)$ from numerals whose
job is alignment or unfinished closure. Readers evaluating the
programme should use the roadmap's honesty table, not footnotes
in a survey note.

%============================================================================
\section{Octonions and the Standard-Model bookkeeping layer}
\label{sec:octonion-bookkeeping}

The
\hyperref[lean:sm-embedding]{SM embedding module},
\hyperref[lean:triality-cycle]{Spin(8) triality module}, and
\hyperref[lean:so8-generators]{$\mathfrak{so}(8)$ generators
module} package octonionic $8$-dimensional carriers,
$\mathfrak{so}(8)$ generators, and Standard-Model quantum
numbers. This layer fixes representation labels; the dynamical
masses above come from the shell programme in
Sections~\ref{sec:core}--\ref{sec:quarks}. In parallel, the
active colour triplet is packaged on the $3\times 3$ chart (the
\hyperref[lean:color-chart]{quark colour-carrier gauge
scaffold} and the
\hyperref[lean:su3-chart-closure]{$\mathrm{SU}(3)$ colour
chart-closure module}) and embedded into the octonion-indexed
carrier via the
\hyperref[lean:color-carrier-closure]{strong-colour carrier
closure module}; an optional Lake target supplies generated
simp lemmas for nonzero $f^{abc}$ structure constants without
enlarging the default Lean cone.

\section*{Acknowledgments}
The author thanks the open-science community for the public
software stack on which the HQIV-LEAN library is built (Lean~4,
\texttt{mathlib}, and the Zenodo HQIV community
archive).\footnote{AI-assisted drafting and repository tooling
were used during manuscript preparation; all formal claims
refer to the HQIV-LEAN library, and the author reviewed
hypothesis lists in the cited Lean modules.}

\appendix

\section{Lean module index}
\label{app:lean-index}

The Lean development lives in the public slimmed mirror
\url{https://github.com/HQIV/hqiv-lean} \cite{hqiv-lean}.
Every reader-friendly anchor used in the body resolves to a
single entry below, with a one-line description and clickable
links to the Lean module and identifiers.  All theorems cited
in the body are sorry-free at the time of writing.

\begingroup
\sloppy
\setlength{\emergencystretch}{4em}
\begin{description}[leftmargin=0em,style=nextline]
\item[\phantomsection\label{lean:auxiliary-T}\textbf{Auxiliary-field temperature ladder.}]
  $T(m)=1/(m+1)$, $\varphi(m)=2(m+1)$, and shell-shape constants.
  Lean: \leanid{T}, \leanid{varphi},
  \leanid{shell\_shape} in
  \leanfile{Hqiv/Geometry/AuxiliaryField.lean}.
\item[\phantomsection\label{lean:phase-lift}\textbf{Phase-lift generator.}]
  Fixed antisymmetric $8\times 8$ generator $\Delta$ with
  nonzero $(1,7),(7,1)$ entries. Lean:
  \leanid{phaseLiftDelta},
  \leanid{phaseLiftCoeff} in
  \leanfile{Hqiv/Algebra/PhaseLiftDelta.lean}.
\item[\phantomsection\label{lean:phase-lift-coeff}\textbf{Phase-lift coefficient.}]
  $\varphi(m)/6$ used in projected mode strengths. Lean:
  \leanid{phaseLiftCoeff} in
  \leanfile{Hqiv/Algebra/PhaseLiftDelta.lean}.
\item[\phantomsection\label{lean:phase-lift-quadratic}\textbf{Phase-lift quadratic-trace identity.}]
  $\operatorname{tr}(\Delta\Delta^{\mathsf T})=2$. Lean:
  \leanid{phaseLiftDeltaQuadraticTrace\_eq\_two} in
  \leanfile{Hqiv/Algebra/PhaseLiftDelta.lean}.
\item[\phantomsection\label{lean:fano-line-trace}\textbf{Fano-line selector-trace identity.}]
  $\operatorname{tr}(\Pi_L)=3$ for the diagonal selector
  associated to a Fano line $L$. Lean:
  \leanid{fanoLineSelector\_trace} in
  \leanfile{Hqiv/Algebra/FanoLine.lean}.
\item[\phantomsection\label{lean:fano-spectrum}\textbf{Fano O-Maxwell spectrum module.}]
  Matrix-level construction of the spectral 1-jet underlying
  the detuning denominator. Lean:
  \leanfile{Hqiv/Physics/FanoOmaxwellSpectrum.lean}.
\item[\phantomsection\label{lean:fano-spectral-norm}\textbf{Spectral projection-normalisation theorem.}]
  $\mathcal{N}(L)=(3+2)/5=1$. Lean:
  \leanid{spectralProjectionNormalization\_eq\_one} in
  \leanfile{Hqiv/Physics/FanoOmaxwellSpectrum.lean}.
\item[\phantomsection\label{lean:fano-rindler-1jet}\textbf{First-order Fano--Rindler jet identity.}]
  $D_L(m)=1+(\gamma/2)\,m$ on the spectral 1-jet. Lean:
  \leanid{spectralFanoRindler1Jet\_eq\_rindler} and
  \leanid{spectralFanoRindler1Jet\_eq\_one\_plus\_half\_gamma}
  in
  \leanfile{Hqiv/Physics/FanoOmaxwellSpectrum.lean}.
\item[\phantomsection\label{lean:rindler-shared}\textbf{Shared Rindler detuning.}]
  The downstream-public detuning $1+\gamma m/2$. Lean:
  \leanid{rindlerDetuningShared} in
  \leanfile{Hqiv/Physics/FanoResonance.lean}.
\item[\phantomsection\label{lean:triality-denominator}\textbf{Triality-projected denominator.}]
  Public detuning denominator coinciding with the spectral
  1-jet. Lean:
  \leanid{trialityProjectedDenominator},
  \leanid{trialityProjectedDenominator\_firstOrder} in
  \leanfile{Hqiv/Physics/FanoTrialityDetuningScaffold.lean}.
\item[\phantomsection\label{lean:triality-denominator-firstorder}\textbf{Triality-denominator first-order corollary.}]
  Same Lean reference as above; the proved corollary status
  is on the \leanid{trialityProjectedDenominator\_firstOrder}
  theorem.
\item[\phantomsection\label{lean:triality-cycle}\textbf{Spin(8) triality module.}]
  Order-3 representation cycle on $(8_v,8_{s^+},8_{s^-})$.
  Lean: \leanfile{Hqiv/Algebra/Triality.lean}.
\item[\phantomsection\label{lean:fano-resonance}\textbf{Fano resonance module.}]
  Shared shell surface $S(m)=(m+1)(m+2)$ and detuned
  outer-shell surface. Lean:
  \leanid{shellSurface},
  \leanid{detunedShellSurface} in
  \leanfile{Hqiv/Physics/FanoResonance.lean}.
\item[\phantomsection\label{lean:detuned-shell}\textbf{Detuned outer-shell surface.}]
  Same Lean reference as above; the
  \leanid{detunedShellSurface} definition is the entry point.
\item[\phantomsection\label{lean:resonance-step}\textbf{Geometric-resonance-step ratio.}]
  Inter-generation step built from detuned outer-shell
  surfaces. Lean:
  \leanid{geometricResonanceStep},
  \leanid{upResonanceProduct},
  \leanid{downResonanceProduct} in
  \leanfile{Hqiv/Physics/QuarkMetaResonance.lean}.
\item[\phantomsection\label{lean:resonance-product}\textbf{Resonance-product factor.}]
  Triality-keyed resonance product used by the geometric SM
  mass functional. Lean:
  \leanid{resonanceProduct} in
  \leanfile{Hqiv/Physics/ChargedLeptonResonance.lean}.
\item[\phantomsection\label{lean:lepton-resonance}\textbf{Charged-lepton resonance module.}]
  Phenomenology layer for $e/\mu/\tau$ ratios. Lean:
  \leanfile{Hqiv/Physics/ChargedLeptonResonance.lean}.
\item[\phantomsection\label{lean:lepton-global-detuning}\textbf{Global-detuning lepton module.}]
  Lightweight charged-lepton ratios on the global detuning
  spine. Lean:
  \leanfile{Hqiv/Physics/LeptonResonanceGlobalDetuning.lean}.
\item[\phantomsection\label{lean:modal-horizon}\textbf{Modal-frequency horizon module.}]
  Primary input layer treating modal frequency / horizon data
  as the lock-in coordinate. Lean:
  \leanfile{Hqiv/Physics/ModalFrequencyHorizon.lean}.
\item[\phantomsection\label{lean:gauge-lepton}\textbf{Derived gauge-and-lepton sector module.}]
  Closed rational witnesses for $M_W,M_Z,m_H$ from the
  geometric vacuum, plus the neutrino ladder. Lean:
  \leanid{vacuumExpectationValueGauge},
  \leanid{boson\_witness\_M\_W},
  \leanid{boson\_witness\_M\_Z},
  \leanid{boson\_witness\_m\_H},
  \leanid{neutrino\_masses\_from\_outer\_horizon} in
  \leanfile{Hqiv/Physics/DerivedGaugeAndLeptonSector.lean}.
\item[\phantomsection\label{lean:gauge-couplings}\textbf{Triality-fixed gauge couplings.}]
  $\mathrm{SU}(2)_w$ and $\mathrm{U}(1)_Y$ couplings from
  triality. Lean:
  \leanid{su2CouplingDerived},
  \leanid{u1CouplingDerived} in
  \leanfile{Hqiv/Physics/DerivedGaugeAndLeptonSector.lean}.
\item[\phantomsection\label{lean:boson-witnesses}\textbf{Rational-boson witness pack.}]
  $M_W^{\mathrm{derived}}=392/5$ and the matching $M_Z,m_H$
  rational identities. Lean:
  \leanid{boson\_witness\_M\_W},
  \leanid{boson\_witness\_M\_Z},
  \leanid{boson\_witness\_m\_H} in
  \leanfile{Hqiv/Physics/DerivedGaugeAndLeptonSector.lean}.
\item[\phantomsection\label{lean:neutrino-ladder}\textbf{Outer-horizon neutrino mass identity.}]
  $m_{\nu_e}^{\mathrm{derived}}=(\gamma/S(M+2))\,M_Z^{\mathrm{derived}}$
  with the generation cascade. Lean:
  \leanid{m\_nu\_e\_derived},
  \leanid{neutrino\_masses\_from\_outer\_horizon} in
  \leanfile{Hqiv/Physics/DerivedGaugeAndLeptonSector.lean}.
\item[\phantomsection\label{lean:weak-higgs}\textbf{Weak/Higgs O-Maxwell scaffold.}]
  Typed extension surface on the proved O-Maxwell action of
  \cite{hqiv-oct-action-paper}, with non-abelian-form $F$,
  Higgs potential, and lock-in vev proxy. Lean:
  \leanfile{Hqiv/Physics/WeakHiggsFromOMaxwellScaffold.lean}.
\item[\phantomsection\label{lean:weak-higgs-reductions}\textbf{Weak/Higgs reduction-lemma pack.}]
  Reduction theorems collapsing the weak/Higgs layer back to
  the abelian O-Maxwell core. Lean:
  \leanid{weakF\_reduces\_to\_abelian},
  \leanid{weakF\_diag\_eq\_zero\_of\_abelian},
  \leanid{lockinVev\_sq\_eq\_eta\_paper},
  \leanid{lockinVev\_eq\_sqrt\_eta\_paper},
  \leanid{action\_O\_weak\_higgs\_reduces\_exactly\_to\_action\_O\_Maxwell},
  \leanid{weakHiggs\_triality\_tilt\_average\_eq\_one} in
  \leanfile{Hqiv/Physics/WeakHiggsFromOMaxwellScaffold.lean}.
\item[\phantomsection\label{lean:lockin-vev}\textbf{Lock-in vev identity.}]
  $v_{\mathrm{lockin}}=\sqrt{\eta_{\mathrm{paper}}\,\Omega_k(M;M)}$.
  Lean: \leanid{lockinVev\_eq\_sqrt\_eta\_paper} in
  \leanfile{Hqiv/Physics/WeakHiggsFromOMaxwellScaffold.lean}.
\item[\phantomsection\label{lean:action}\textbf{O-Maxwell action module.}]
  The proved abelian O-Maxwell action whose
  Euler--Lagrange identity is the upstream variational base
  \cite{hqiv-oct-action-paper}. Lean:
  \leanfile{Hqiv/Physics/Action.lean}.
\item[\phantomsection\label{lean:alpha-ratio}\textbf{Lattice $\alpha$-ratio identity.}]
  $\alpha=3/5$ at every truncation, imported from
  \cite{hqiv-lightcone-paper,hqiv-so8-paper}. Lean:
  \leanid{latticeAlphaRatio\_eq\_alpha} in
  \leanfile{Hqiv/Geometry/OctonionicLightCone.lean}.
\item[\phantomsection\label{lean:content-mass-bridge}\textbf{Conserved-content mass-scaling bridge.}]
  Rank-$l$ scaling on top of the shell imprint. Lean:
  \leanid{massScalingAnsatz} in
  \leanfile{Hqiv/Physics/ConservedContentMassBridge.lean}.
\item[\phantomsection\label{lean:scale-witness}\textbf{Scale-witness enum.}]
  Inductive type with the proton lock-in / CODATA-$\alpha$ /
  present-CMB modes. Lean:
  \leanid{ScaleWitness} in
  \leanfile{Hqiv/Physics/ScaleWitness.lean}.
\item[\phantomsection\label{lean:proton-lockin}\textbf{Proton lock-in scale-witness.}]
  Default mode pinning the mass/unit chart at $M$. Lean:
  \leanid{ScaleWitness.proton\_lockin} in
  \leanfile{Hqiv/Physics/ScaleWitness.lean}.
\item[\phantomsection\label{lean:hqiv-witnesses}\textbf{HQIV witnesses Lake target.}]
  The Lake target that exports the closure-witness JSON
  bundle. Source roots are summarised in the repository
  README; the target is invoked as
  \texttt{lake build HQIVWitnesses}.
\item[\phantomsection\label{lean:witnesses-json}\textbf{HQIV witness JSON bundle.}]
  Closure witnesses ($W,Z,H$ rational identities, neutrino
  ladder, proton anchor, etc.) used by the downstream
  pipeline. File: \nolinkurl{data/hqiv_witnesses.json}.
\item[\phantomsection\label{lean:info-energy}\textbf{Informational-energy module.}]
  $E_{\mathrm{tot}}=m_{\mathrm{rest}}+1/\Theta_{\mathrm{local}}(\xi)$
  with additive vs.\ multiplicative readout gauges. Lean:
  \leanfile{Hqiv/Physics/InformationalEnergyMass.lean}.
\item[\phantomsection\label{lean:additive-loc}\textbf{Additive-localisation readout gauge.}]
  Boson convention: full $E_{\mathrm{tot}}$ is the
  observable mass. Lean:
  \leanid{additiveLocalization} in
  \leanfile{Hqiv/Physics/InformationalEnergyMass.lean}.
\item[\phantomsection\label{lean:lapse-readout}\textbf{Multiplicative-lapse readout gauge.}]
  Hadron convention: rest slot $\div N_{\mathrm{lapse}}$
  with localisation in the lapse increment. Lean:
  \leanid{multiplicativeLapse},
  \leanid{hadronMassFromXi} in
  \leanfile{Hqiv/Physics/InformationalEnergyMass.lean};
  lapse-readout lemmas live in
  \leanfile{Hqiv/Physics/LapseMassReadout.lean}.
\item[\phantomsection\label{lean:gauge-transform}\textbf{Gauge-transformation lemma.}]
  When the additive and multiplicative readout gauges agree
  after $m_{\mathrm{rest}}$ calibration. Lean:
  \leanid{gauge\_transformation\_localization\_to\_lapse} in
  \leanfile{Hqiv/Physics/InformationalEnergyMass.lean}.
\item[\phantomsection\label{lean:cont-xi-coupling}\textbf{Continuous-$\xi$ coupling module.}]
  Fano coupling system rows in Lean form. Lean:
  \leanfile{Hqiv/Physics/ContinuousXiCoupling.lean}.
\item[\phantomsection\label{lean:info-energy-row}\textbf{Informational-energy mass-row identity.}]
  $c_0+\mathrm{loc}(\xi_G)=2\pi\,\Omega_k(\xi_G)$ budget
  identity. Lean:
  \leanid{informationalEnergyMassRow},
  \leanid{informationalEnergyMassRow\_budget} in
  \leanfile{Hqiv/Physics/InformationalEnergyMass.lean}.
\item[\phantomsection\label{lean:info-energy-row-twopi}\textbf{$2\pi\,\Omega_k$ specialisation.}]
  $E_{\mathrm{tot}}/(2\pi)=\Omega_k$ when the row holds.
  Lean:
  \leanid{informationalEnergy\_satisfied\_when\_row\_holds},
  \leanid{informationalEnergy\_over\_twoPi\_eq\_omegaK\_when\_row\_holds}
  in
  \leanfile{Hqiv/Physics/InformationalEnergyMass.lean}.
\item[\phantomsection\label{lean:meta-horizon}\textbf{Meta-horizon excited-states module.}]
  Baryon excitations as radial $(n)$ and orbital $(\ell)$
  modes on the lock-in drum. Lean:
  \leanid{totalModeMass} in
  \leanfile{Hqiv/Physics/MetaHorizonExcitedStates.lean}.
\item[\phantomsection\label{lean:baryogenesis}\textbf{Baryogenesis module.}]
  Records $m_{\mathrm{lockin}}=M=\mathrm{referenceM}$ and
  related lock-in identities. Lean:
  \leanfile{Hqiv/Physics/Baryogenesis.lean}.
\item[\phantomsection\label{lean:bound-states}\textbf{Bound-states module.}]
  Shell-resolved $\alpha_{\mathrm{eff}}(m)$ and binding
  bookkeeping. Lean:
  \leanfile{Hqiv/Physics/BoundStates.lean}.
\item[\phantomsection\label{lean:harmonic-ladder}\textbf{Harmonic-ladder mass module.}]
  Hydrogenic $\alpha^2$ scalings on the harmonic ladder.
  Lean:
  \leanfile{Hqiv/Physics/HarmonicLadderMass.lean}.
\item[\phantomsection\label{lean:quark-meta}\textbf{Quark meta-resonance module.}]
  Quark ladder: GeV anchor (top), $\mathbb{N}$ shell
  triples, geometric-resonance steps. Lean:
  \leanfile{Hqiv/Physics/QuarkMetaResonance.lean}.
\item[\phantomsection\label{lean:t-lockin-id}\textbf{Lock-in temperature identity.}]
  $T_{\mathrm{lockin}}=T(M)$ as the lock-in horizon
  temperature. Lean:
  \leanid{T\_lockin\_eq\_temperature\_at\_quark\_top\_birth\_shell}
  in
  \leanfile{Hqiv/Physics/SM_GR_Unification.lean}.
\item[\phantomsection\label{lean:color-resonance}\textbf{Color-resonance mass ladder.}]
  Heavy / mid / light bands of the colour-composed
  visible-state API. Lean:
  \leanid{allowedColorResonanceMass},
  \leanid{topLockinColorResonanceAnchorMass} in
  \leanfile{Hqiv/Physics/QuarkMetaResonance.lean}.
\item[\phantomsection\label{lean:derived-nucleons}\textbf{Derived nucleon-mass module.}]
  Proton, neutron, and neutron--proton-gap definitions and
  positivity. Lean:
  \leanid{derivedProtonMass},
  \leanid{derivedNeutronMass},
  \leanid{derivedDeltaM} in
  \leanfile{Hqiv/Physics/DerivedNucleonMass.lean}.
\item[\phantomsection\label{lean:derived-proton}\textbf{Derived proton mass.}]
  Same Lean reference as above (\leanid{derivedProtonMass}).
\item[\phantomsection\label{lean:delta-m}\textbf{Neutron--proton gap identity.}]
  Same Lean reference as above (\leanid{derivedDeltaM}).
\item[\phantomsection\label{lean:shared-binding}\textbf{Shared-binding nucleon energy.}]
  Composite-trace binding scale at $M$ used by both
  ground-state nucleons. Lean:
  \leanid{nucleonSharedBinding\_MeV} in
  \leanfile{Hqiv/Physics/DerivedNucleonMass.lean}.
\item[\phantomsection\label{lean:proton-anchor}\textbf{Proton-anchor mass constant.}]
  $M_p^{\mathrm{anchor}}=938.272\,\mathrm{MeV}$ kept for
  legacy lapse-readout lemmas. Lean:
  \leanid{protonAnchorMass\_MeV} in
  \leanfile{Hqiv/Physics/QuarkMetaResonance.lean}.
\item[\phantomsection\label{lean:qcd-shell}\textbf{QCD-onset shell pin.}]
  $M=\mathrm{referenceM}$ as the QCD-onset substrate pin.
  Lean: \leanid{qcdShell},
  \leanid{latticeStepCount} in
  \leanfile{Hqiv/Geometry/AuxiliaryField.lean}.
\item[\phantomsection\label{lean:sm-gr}\textbf{SM--GR unification module.}]
  $\alpha_{\mathrm{GUT}}=1/42$ from cube directions $\times$
  octonion imaginary units; geometric SM mass functional;
  lock-in temperature identity. Lean:
  \leanfile{Hqiv/Physics/SM_GR_Unification.lean}.
\item[\phantomsection\label{lean:sm-mass}\textbf{Geometric SM mass functional.}]
  $\tau$-Planck scale divided by the resonance product. Lean:
  \leanid{smMassFromGeometry},
  \leanid{smMassFromGeometryLabel},
  \leanid{m\_top\_quark\_natural} in
  \leanfile{Hqiv/Physics/SM_GR_Unification.lean}.
\item[\phantomsection\label{lean:sm-embedding}\textbf{SM embedding module.}]
  Octonionic 8-dimensional carriers and Standard-Model
  representation labels. Lean:
  \leanfile{Hqiv/Algebra/SMEmbedding.lean}.
\item[\phantomsection\label{lean:so8-generators}\textbf{$\mathfrak{so}(8)$ generators module.}]
  Generator bookkeeping for the bookkeeping layer. Lean:
  \leanfile{Hqiv/Algebra/Generators.lean}.
\item[\phantomsection\label{lean:color-chart}\textbf{Quark colour-carrier gauge scaffold.}]
  Abstract $3$-component colour chart. Lean:
  \leanfile{Hqiv/Physics/QuarkColorCarrierGaugeScaffold.lean}.
\item[\phantomsection\label{lean:su3-chart-closure}\textbf{$\mathrm{SU}(3)$ colour chart-closure module.}]
  Hermitian Gell-Mann data, half-generators, and the real
  totally antisymmetric $f^{abc}$ tensor. Lean:
  \leanfile{Hqiv/Physics/StrongColorSu3ChartClosure.lean}.
\item[\phantomsection\label{lean:color-carrier-closure}\textbf{Strong-colour carrier-closure module.}]
  Embedding of the $3\times 3$ chart into the
  octonion-indexed
  $\mathrm{EuclideanSpace}\,\mathbb{C}\,(\mathrm{Fin}\,8)$
  carrier shared with the electroweak layer. Lean:
  \leanfile{Hqiv/Physics/StrongColorCarrierClosure.lean}.
\end{description}
\endgroup

\bibliographystyle{plain}
\bibliography{../../references}

\end{document}
